\section{Conclusion}
\label{sec:conclusion}

We presented three configurable stream mapping strategies for running MQTT over QUIC. Through six experiments on controlled infrastructure, we demonstrated that per-topic QUIC stream mapping provides measurable HOL blocking isolation: inter-topic latency spread amplifies 192$\times$ under 5\% packet loss, confirming that independent stream loss recovery allows topics to experience different retransmission delays rather than stalling together as TCP does (1.3$\times$ spread amplification). Spike isolation ratio corroborates this---at 1\% loss, 28\% of per-topic QUIC latency spikes affect individual topics independently, versus only 1\% for TCP. The isolation is partial rather than complete due to QUIC's per-connection congestion control, which creates correlated congestion responses across streams.

Our evaluation reveals a nuanced throughput tradeoff: at baseline (0\% loss), TCP achieves 1.5--2.4$\times$ the throughput of all QUIC strategies, but under loss ($\geq$1\%) QUIC outperforms TCP by up to 2.7$\times$ due to more efficient loss recovery. Per-topic mapping offers throughput within 13--32\% of single-stream QUIC depending on connection count. QUIC datagrams achieve 19\% lower latency than streams at 5\% packet loss with comparable throughput, making them suitable for loss-tolerant telemetry.

These results provide practitioners with concrete guidance for deploying MQTT over QUIC in IoT systems: use per-topic mapping for multi-topic workloads on lossy networks where independent stream recovery provides meaningful isolation and QUIC's throughput advantage under loss is beneficial, and use TCP when baseline throughput on reliable networks is the dominant concern. The implementation is available as open-source software.
